\chapter{Hemoglobin Electrophoresis}

\section*{What Is This Test?}
Hemoglobin electrophoresis is a laboratory technique that separates the different types of hemoglobin in a blood sample by applying an electric field. Hemoglobin (Hb) is the oxygen-carrying protein inside red blood cells, composed of two alpha-like globin chains and two beta-like globin chains assembled around a heme group \cite{steinberg2009disorders}. In a healthy adult, the major hemoglobin is HbA (alpha-2 beta-2, approximately 95--98\%), with small amounts of HbA2 (alpha-2 delta-2, 1.5--3.5\%) and HbF (alpha-2 gamma-2, less than 2\%). Hemoglobin electrophoresis uses cellulose acetate or agarose gel at alkaline pH to separate hemoglobins by their net electrical charge and migration distance, producing characteristic bands: HbA migrates fastest, followed by HbF, then HbA2, then HbS, and finally HbC. Because each hemoglobin variant has a characteristic migration pattern, electrophoresis can identify hemoglobinopathies including sickle cell disease (HbSS), sickle cell trait (HbAS), hemoglobin C disease (HbCC), hemoglobin SC disease (HbSC), and the thalassemias. High-performance liquid chromatography (HPLC) and capillary electrophoresis have largely replaced gel electrophoresis in modern laboratories because they are fully automated, quantitative, and more sensitive for detecting low-abundance variants, but the interpretive principles are the same.

\section*{Alternative Names}
\begin{itemize}
  \item Hb Electrophoresis
  \item Hemoglobin Fractionation
  \item Hemoglobinopathy Screen
  \item Sickle Cell Screen (qualitative)
  \item Hemoglobin HPLC
  \item Hemoglobin Capillary Electrophoresis
\end{itemize}
\section*{Principle}
Electrophoretic separation exploits differences in the net charge of hemoglobin molecules at a given pH. At alkaline pH (8.4--8.6), hemoglobins are negatively charged and migrate toward the anode in an electric field. The rate of migration depends on amino acid substitutions that alter surface charge. HbA migrates farthest; HbF and HbA2 migrate more slowly; HbS migrates between HbA2 and HbA; and HbC migrates slowest. After separation, the gel is stained and the bands are identified by their position relative to known controls. Quantification is achieved by densitometry (measuring the optical density of each stained band) or, on modern platforms, by integration of HPLC elution peaks or capillary electrophoretic traces. HPLC additionally exploits differences in hydrophobicity, separating hemoglobins as they elute from a cation-exchange column; retention times are compared against a standard library of common variants. Because HbS (glutamate-to-valine substitution at beta-6) and HbC (glutamate-to-lysine at beta-6) both alter charge, they are readily separated by both methods. Confirmation of an abnormal band is typically performed with a complementary method, such as citrate agar electrophoresis at acid pH, isoelectric focusing, or DNA-based (molecular) testing.

\section*{History}
Hemoglobin was first separated electrophoretically in the 1950s when Linus Pauling and colleagues demonstrated that sickle hemoglobin migrates differently from normal hemoglobin on an electric field, coining the term ``molecular disease'' for sickle cell anemia. This discovery, published in 1949, was the first demonstration that a specific molecular alteration underlies a human disease \cite{pauling1949sickle}. Starch gel and then cellulose acetate electrophoresis became standard clinical tools in the 1960s. Janet Watson's earlier observation of the benign infant course of sickle cell anemia had hinted at the protective role of HbF, which became measurable with the development of alkali denaturation and later electrophoretic techniques. In the 1970s--1980s, isoelectric focusing improved resolution of variants with identical charge. High-performance liquid chromatography, introduced for hemoglobin analysis in the 1980s, became the reference method for quantification of HbA2 and HbF and for newborn screening for hemoglobinopathies. Capillary electrophoresis emerged in the 2000s. Newborn screening programs in many countries now use HPLC or isoelectric focusing to detect hemoglobinopathies on dried blood spots within the first days of life \cite{michlitsch2009newborn}.

\section*{Why This Test Is Ordered}
\begin{itemize}
  \item Diagnosis and classification of sickle cell disease: hemoglobin electrophoresis confirms the presence of HbS and determines whether the patient has sickle cell trait (HbAS), sickle cell anemia (HbSS), or compound heterozygosity (e.g., HbSC disease), which differ markedly in clinical severity
  \item Evaluation of hemolytic anemia: a positive family history, chronic hemolysis, or anemia of unknown cause may prompt screening for hemoglobinopathy
  \item Diagnosis of the thalassemias: elevated HbA2 (3.5--7\%) with or without elevated HbF supports the diagnosis of beta-thalassemia trait; reduced or absent beta-chain production with relative excess of HbA2 and HbF characterizes beta-thalassemia major and intermedia
  \item Newborn screening: mandated in many jurisdictions to detect sickle cell disease and severe thalassemia before symptoms develop, enabling penicillin prophylaxis and pneumococcal vaccination
  \item Confirmatory testing after an abnormal newborn screen or a positive sickle cell solubility test (sickle prep), which can only detect HbS, not distinguish trait from disease \cite{ryan2010significant}
  \item Prenatal diagnosis and genetic counseling in couples at risk for hemoglobinopathies, including those of African, Mediterranean, Middle Eastern, Southeast Asian, and Indian descent
  \item Preoperative and preconception screening in at-risk populations to identify carriers who may transmit disease to offspring
  \item Assessment of unexplained microcytic anemia in populations where hemoglobinopathy and iron deficiency coexist
\end{itemize}

\section*{Primary vs Secondary Test}
This is a secondary confirmatory test. Screening is initially performed with a CBC and red blood cell indices (microcytic anemia, abnormal RDW) and, when sickle hemoglobin is suspected, a solubility test. Hemoglobin electrophoresis (or HPLC/capillary electrophoresis) is the definitive test that identifies and quantifies the specific hemoglobin variant. In newborn screening programs, however, HPLC or isoelectric focusing is used as the primary analytical method on dried blood spots.

\section*{How This Test Is Performed}
A blood sample is collected by standard venipuncture into an EDTA (purple-top) tube. A CBC with red cell indices is usually obtained simultaneously. In the laboratory, red cells are washed and lysed to release hemoglobin. The hemolysate is applied to the separation medium and an electric current is applied. On HPLC platforms, the hemolysate is injected directly onto the analytical column. After separation, results are reported as a quantitative percentage of each hemoglobin fraction (HbA, HbA2, HbF) and any variant bands are identified by name. For newborn screening, a heel-stick sample is applied to a filter paper card (dried blood spot) that is analyzed directly. Results are typically available within 1--3 days for routine testing.

\section*{What Preparation Is Needed}
No fasting or special preparation is required. Recent blood transfusion within the preceding 8--12 weeks can obscure the diagnosis because transfused donor cells contribute normal hemoglobin; if transfusion is unavoidable, a sample obtained before transfusion is preferred for diagnostic accuracy. In patients with recent transfusion, molecular testing of a pre-transfusion sample or buccal cells may be required.

\section*{Contraindications and Precautions}
\begin{itemize}
  \item There are no absolute contraindications to standard venipuncture. Important interpretive cautions:
  \item (1) Transfusion within the past 3 months dilutes the patient's own hemoglobin, potentially making a variant (especially HbS in sickle cell trait or HbC trait) appear at artificially low percentages or be missed entirely.
  \item (2) In newborns, HbF predominates, and definitive diagnosis of beta-thalassemia may require repeat testing after 6--12 months when HbF declines.
  \item (3) Severe iron deficiency can lower HbA2 levels, masking the elevated HbA2 of beta-thalassemia trait.
  \item (4) Patients on hydroxyurea therapy have elevated HbF, which can complicate interpretation of HbF-dependent patterns.
  \item (5) Rapidly degraded specimens, hemoglobin degradation products, and high levels of HbH or Hb Bart's in alpha-thalassemia can produce atypical patterns. Any result that is inconsistent with the clinical picture should be confirmed by a second method or by DNA analysis.
\end{itemize}
\section*{Risks}
Minimal risk. Slight pain or bruising at the venipuncture site. Rarely: hematoma, infection, or vasovagal syncope.

\section*{What Happens After the Test}
Results are available within 1--3 days for routine samples and are reported as the percentage distribution of hemoglobin fractions plus the identification of any variants. The results are interpreted in the context of the CBC, red cell indices, iron studies, and clinical history. A definitive diagnosis of a clinically significant hemoglobinopathy is followed by genetic counseling, which is especially important for newborns, reproductive-age carriers, and couples planning a pregnancy. For patients with sickle cell disease, established follow-up includes newborn penicillin prophylaxis, pneumococcal and meningococcal vaccination, transcranial Doppler screening for stroke risk, annual ophthalmologic and renal screening, and referral to a hematologist.

\section*{Understanding Results}

\subsection*{Below Normal (Absence or Reduction of a Normal Fraction)}
\begin{itemize}
  \item \textbf{Elevated HbA2 with normal or near-normal HbA (trait):} HbA2 3.5--7\% with microcytic anemia and normal or slightly reduced HbA is diagnostic of beta-thalassemia trait (heterozygous beta-thalassemia) \cite{thein2013molecular}. The carrier is usually asymptomatic but can transmit beta-thalassemia major to offspring.
  \item \textbf{Reduced or absent HbA with high HbF (homozygous beta-thalassemia):} In beta-thalassemia major, HbA is absent or minimal, HbF is markedly elevated (30--98\%), and HbA2 is variable. Clinically severe transfusion-dependent anemia beginning in infancy.
  \item \textbf{High HbF with normal HbA2 in an adult:} Hereditary persistence of fetal hemoglobin, delta-beta thalassemia, aplastic anemia, or myeloproliferative disorders; on hydroxyurea therapy, HbF may rise substantially as a desired therapeutic effect.
  \item \textbf{Reduced HbA2 and HbA with fast-migrating bands (alpha-thalassemia):} Alpha-thalassemia trait may show a mild reduction of HbA2 (below 2.5\%); HbH (beta-4) or Hb Bart's (gamma-4) appear as fast-migrating bands in the more severe forms.
\end{itemize}

\subsection*{Above Normal (Presence of a Variant Hemoglobin)}
\begin{itemize}
  \item \textbf{HbAS (sickle cell trait):} HbA 55--65\%, HbS 35--45\%, normal HbA2 and HbF. Generally benign; individuals have normal life expectancy but inherit a 50\% risk of transmitting the HbS gene.
  \item \textbf{HbSS (sickle cell anemia):} HbS 80--95\%, HbF 2--20\%, HbA2 mildly elevated, no HbA. Causes chronic hemolysis, vaso-occlusive crises, acute chest syndrome, stroke, and organ damage \cite{bunn1997pathogenesis}.
  \item \textbf{HbSC disease:} HbS 50\%, HbC 50\% approximately, with elevated HbF. Similar to but often milder than HbSS; increased risk of retinopathy and aseptic necrosis.
  \item \textbf{HbCC (hemoglobin C disease):} HbC greater than 90\%, no HbA. Mild hemolytic anemia, splenomegaly, and target cells on smear.
  \item \textbf{HbAC (hemoglobin C trait):} HbA approximately 60--70\%, HbC 30--40\%. Benign carrier state.
  \item \textbf{Other variants (HbE, HbD, HbO-Arab, Hb Lepore):} Characteristic migration patterns on HPLC or capillary electrophoresis; HbE is extremely common in Southeast Asia and, when co-inherited with beta-thalassemia, produces clinically significant disease.
\end{itemize}

\section*{Normal Results}
\begin{itemize}
  \item \textbf{Adult HbA:} 95--98\% of total hemoglobin
  \item \textbf{Adult HbA2:} 1.5--3.5\% (upper limit of normal varies by laboratory; greater than 3.5--4.0\% is suspicious for beta-thalassemia trait)
  \item \textbf{Adult HbF:} less than 2\% (may be up to 5\% in some reference populations)
  \item \textbf{Newborn HbF:} 60--90\% at birth, declining to near-adult levels by 6--12 months
  \item \textbf{Newborn HbA:} 10--40\% at birth
  \item No abnormal hemoglobin variants detected on the standard screen
\end{itemize}

\section*{Follow-up Tests}
\begin{itemize}
  \item \textbf{Complete blood count with red cell indices:} Essential companion testing; microcytosis, hypochromia, and anisocytosis suggest thalassemia, while normal indices with a variant suggest a structural hemoglobinopathy
  \item \textbf{Iron studies (ferritin, iron, TIBC):} Exclude coexistent iron deficiency, which can lower HbA2 and obscure beta-thalassemia trait
  \item \textbf{DNA-based molecular testing:} Definitive identification of the specific mutation; essential for prenatal diagnosis, evaluation of complex or rare variants, and cases where transfusion confounds protein-based testing
  \item \textbf{Reticulocyte count and peripheral smear:} Assess the degree of hemolysis and look for characteristic findings (target cells in HbC and HbSC disease, sickle cells in HbSS)
  \item \textbf{Bilirubin, LDH, haptoglobin:} Laboratory markers of hemolysis
  \item \textbf{Genetic counseling and partner testing:} For carriers and affected individuals to define reproductive risk
  \item \textbf{HbF quantification (Kleihauer-Betke or HPLC):} When hereditary persistence of fetal hemoglobin or hydroxyurea response is being evaluated
\end{itemize}

\section*{References}
\bibliographystyle{unsrtnat}
\bibliography{references}

\clearpage
\section*{Regulatory Status and Authorization Date}
This chapter describes a test category, not a specific commercial product. FDA, CDSCO, EU IVDR, and MHRA approval or registration dates therefore cannot be assigned without the assay manufacturer, product name, instrument, and intended use. For laboratory-developed tests, an individual agency approval date may not exist. See the Preface for the interpretation of regulatory status.

\clearpage
